
\section{Pendências e riscos conhecidos}

\begin{enumerate}
    \item Importar \texttt{SEW} em \texttt{model/executor.py} e repetir
    toda a suíte de testes.

    \item Validar explicitamente o domínio de \texttt{requested\_vl},
    \texttt{sew}, \texttt{lmul}, índices de registradores e tamanhos de
    máscaras.

    \item Definir se as instruções de memória terão EEW fixo no mnemônico,
    como \texttt{vle8.v}, \texttt{vle16.v} e \texttt{vle32.v}, ou se
    haverá uma operação genérica dependente do SEW ativo. A decisão deve
    evitar que o nome da instrução contradiga seu comportamento.

    \item Generalizar o gerador strip-mined para receber SEW e LMUL e
    calcular VLMAX, tamanho do elemento e alinhamento de registradores
    de forma dinâmica.

    \item Revisar a representação de \texttt{mem\_writes} no trace para
    que seja consistente quanto ao tipo das chaves de endereço e ao
    formato dos valores.

    \item Formalizar a semântica de troca de SEW e LMUL entre instruções,
    em especial como as políticas de cauda se relacionam com a
    reinterpretação física dos registradores.

    \item Especificar estados de erro e acessos inválidos à memória.
    O modelo atual usa exceções Python, mas ainda não define equivalentes
    arquiteturais de traps.

    \item Distinguir claramente EEW de memória e SEW. A simplificação
    atual pode tratar ambos como iguais, mas isso deve ser registrado
    como limitação até que a semântica RVV correspondente seja
    implementada.

    \item Determinar se a meta é uma extensão customizada inspirada em
    RVV ou um subconjunto compatível com uma versão específica da
    especificação RVV. O formato atual de programa Python não é, por si
    só, um encoding binário RISC-V.
\end{enumerate}

\section{Critérios para o próximo marco}

A especificação executável poderá ser considerada estável para esta
etapa quando:

\begin{itemize}
    \item todos os testes passarem a partir da raiz do projeto;
    \item os testes incluírem execução de programas com SEW e LMUL
    explicitamente variados;
    \item o strip-mining funcionar com mais de uma configuração de
    elemento e grupo;
    \item o trace JSON representar corretamente as escritas de memória;
    \item as regras de validação estiverem descritas e testadas;
    \item um ou mais programas de referência tiverem resultados
    reproduzíveis e documentados.
\end{itemize}

\section{Migração futura para SystemC}

A próxima etapa de modelagem deverá preservar a separação conceitual
entre:

\begin{itemize}
    \item estado arquitetural;
    \item interpretação e validação de instruções;
    \item execução funcional;
    \item registro de efeitos observáveis.
\end{itemize}

A primeira versão SystemC poderá continuar sem modelar temporização
detalhada, mas deverá manter o mesmo comportamento arquitetural e
produzir resultados comparáveis com os traces do Python. A modelagem
de ciclos, latências e protocolos deve ser adicionada em uma etapa
separada, para não misturar divergências funcionais com diferenças
temporais.

\section{Histórico deste documento}

\begin{longtable}{p{2.5cm}p{10.5cm}}
\toprule
\textbf{Versão} & \textbf{Descrição} \\
\midrule
0.1 & Registro inicial dos parâmetros, da semântica funcional e do
estado observado do modelo Python. Inclui pendências identificadas na
suíte de testes. \\
\bottomrule
\end{longtable}